\documentclass[11pt]{article}

% ---- PiLens User Guide. British English. Designed to be skimmed. ----
\usepackage[a4paper,margin=20mm,marginparwidth=22mm,marginparsep=5mm]{geometry}
\usepackage[T1]{fontenc}
\usepackage[utf8]{inputenc}
\usepackage{helvet}
\renewcommand{\familydefault}{\sfdefault}
\usepackage{microtype}
\usepackage{graphicx}
\usepackage{xcolor}
\usepackage[most]{tcolorbox}
\usepackage{enumitem}
\usepackage{titlesec}
\usepackage{fancyhdr}
\usepackage{marginnote}
\usepackage{pifont}
\usepackage[hidelinks]{hyperref}
\usepackage{tikz}
\usetikzlibrary{calc}

% ---- palette (matches the appliance UI) ----
\definecolor{ink}{HTML}{11161D}
\definecolor{panel}{HTML}{1B222C}
\definecolor{line}{HTML}{2A323D}
\definecolor{amber}{HTML}{E0922F}
\definecolor{amberlt}{HTML}{F2A93B}
\definecolor{cyan}{HTML}{2C92A4}
\definecolor{cyanlt}{HTML}{5CCFE0}
\definecolor{good}{HTML}{2E9E63}
\definecolor{bad}{HTML}{D4504E}
\definecolor{muted}{HTML}{6A7482}
\definecolor{paper}{HTML}{FBFBF9}

\hypersetup{colorlinks=true, linkcolor=cyan, urlcolor=cyan}
\pagecolor{white}
\color{ink}

% ---- headings ----
\titleformat{\section}{\Large\bfseries\color{ink}}{\thesection}{0.6em}{}[{\color{line}\titlerule[1pt]}]
\titleformat{\subsection}{\large\bfseries\color{amber}}{}{0em}{}
\titleformat{\subsubsection}{\bfseries\color{ink}}{}{0em}{}
\titlespacing*{\section}{0pt}{2.4ex}{1.2ex}
\titlespacing*{\subsection}{0pt}{2ex}{0.6ex}
\setlength{\parindent}{0pt}
\setlength{\parskip}{0.6ex}
\renewcommand{\arraystretch}{1.25}

% ---- callout boxes ----
\newtcolorbox{tldr}{enhanced, colback=amberlt!10, colframe=amberlt, boxrule=0pt,
  leftrule=3pt, arc=2pt, left=8pt, right=8pt, top=5pt, bottom=5pt, breakable,
  fonttitle=\bfseries, title={\ding{220}\ \ TL;DR}}
\newtcolorbox{heads}{enhanced, colback=bad!8, colframe=bad, boxrule=0pt,
  leftrule=3pt, arc=2pt, left=8pt, right=8pt, top=5pt, bottom=5pt, breakable,
  fonttitle=\bfseries, title={\ding{43}\ \ Heads-up}}
\newtcolorbox{tip}{enhanced, colback=cyanlt!10, colframe=cyan, boxrule=0pt,
  leftrule=3pt, arc=2pt, left=8pt, right=8pt, top=5pt, bottom=5pt, breakable,
  fonttitle=\bfseries, title={\ding{42}\ \ Pro tip}}
\newtcolorbox{truth}{enhanced, colback=ink!5, colframe=ink, boxrule=0pt,
  leftrule=3pt, arc=2pt, left=8pt, right=8pt, top=5pt, bottom=5pt, breakable,
  fonttitle=\bfseries, title={\ding{45}\ \ The honest truth}}
\newtcolorbox{shotbox}{enhanced, colback=white, colframe=line, boxrule=0.6pt,
  arc=3pt, left=3pt, right=3pt, top=3pt, bottom=3pt, drop fuzzy shadow=line}

\newcommand{\screen}[2]{\begin{shotbox}\includegraphics[width=\linewidth]{fig/#1}\end{shotbox}\vspace{-0.3ex}{\small\color{muted} Figure: #2}\par\vspace{1ex}}
\newcommand{\key}[1]{\tikz[baseline=(k.base)]{\node[fill=panel,text=white,rounded corners=2pt,inner xsep=4pt,inner ysep=1.5pt,font=\footnotesize\ttfamily](k){#1};}}
\newcommand{\yes}{\textcolor{good}{\ding{51}}}
\newcommand{\no}{\textcolor{bad}{\ding{55}}}
\newcommand{\term}[1]{\textbf{#1}}
\newcommand{\mn}[1]{\marginnote{\footnotesize\color{muted}#1}}

\pagestyle{fancy}\fancyhf{}
\renewcommand{\headrulewidth}{0.4pt}
\lhead{\footnotesize\color{muted}\textbf{PiLens} user guide}
\rhead{\footnotesize\color{muted}\leftmark}
\cfoot{\footnotesize\color{muted}\thepage}
\renewcommand{\sectionmark}[1]{\markboth{#1}{}}

\begin{document}

% ======================= COVER =======================
\begin{titlepage}
\color{ink}
\begin{tikzpicture}[remember picture, overlay]
  \fill[ink] (current page.north west) rectangle ([yshift=-95mm]current page.north east);
  \node[anchor=north west, text=white] at ([xshift=20mm,yshift=-22mm]current page.north west)
    {\fontsize{46}{50}\selectfont\bfseries PiLens};
  \node[anchor=north west, text=amberlt] at ([xshift=21mm,yshift=-38mm]current page.north west)
    {\Large network forensics for the Raspberry Pi};
  % lens mark
  \begin{scope}[shift={([xshift=-42mm,yshift=-42mm]current page.north east)}]
    \draw[amberlt, line width=2.2pt] (0,0) circle (10mm);
    \draw[amberlt, line width=3pt, line cap=round] (7mm,-7mm) -- (15mm,-15mm);
    \draw[cyanlt, line width=1.4pt, line join=round] (-5mm,0) -- (-3mm,0) -- (-1.5mm,4mm) -- (1.5mm,-4mm) -- (3mm,0) -- (5mm,0);
  \end{scope}
\end{tikzpicture}
\vspace{100mm}

{\LARGE\bfseries The friendly field guide}\\[2mm]
{\large\color{muted} Turn a Raspberry Pi into a Wi-Fi test network you can actually \emph{see} inside.}\\[10mm]

\begin{tldr}
This guide is built for skimming. Every section opens with a \textbf{TL;DR}. Boxes
in the margins and colour-coded callouts carry the important bits. \textbf{Read the
box, skip the prose, come back if you need it.} You will not miss anything vital
by jumping around --- that is the point.
\end{tldr}

\vfill
{\footnotesize\color{muted} Version 0.3 \quad\textbullet\quad AGPL-3.0-or-later \quad\textbullet\quad For inspecting devices and networks you are authorised to inspect. Nothing else.}
\end{titlepage}

% ======================= HOW TO READ =======================
\section*{How to read this (30 seconds)}
\addcontentsline{toc}{section}{How to read this}

\begin{tabular}{@{}p{3.2cm}p{12.5cm}@{}}
\ding{220}\ \textbf{TL;DR} & The one-line version. If you read nothing else, read these.\\
\ding{43}\ \textbf{Heads-up} & Something that will bite you if you ignore it.\\
\ding{42}\ \textbf{Pro tip} & A shortcut or nicety. Safe to skip.\\
\ding{45}\ \textbf{Honest truth} & Where PiLens tells you what it \emph{cannot} do. Read these --- they save hours of confusion.\\
\end{tabular}

\vspace{1ex}
\begin{tip}
In a hurry? Jump straight to \textbf{\S1 The ten-minute start}, connect a phone,
and open \textbf{Live Capture}. Come back for the rest when you are curious.
\end{tip}

{\small\color{muted}\tableofcontents}
\newpage

% ======================= 1. QUICK START =======================
\section{The ten-minute start}
\begin{tldr}
Flash the SD card \textgreater\ plug in \textbf{Ethernet} \textgreater\ run one installer
\textgreater\ open the web page \textgreater\ finish the wizard \textgreater\ connect a
test device to the new Wi-Fi \textgreater\ watch its traffic. That is the whole dance.
\end{tldr}

\subsection{What you need}
\begin{itemize}[leftmargin=1.4em, itemsep=1pt]
  \item A Raspberry Pi \textbf{4 or 5} running Raspberry Pi OS Lite (64-bit).
  \item An \textbf{Ethernet} cable to your normal network. \mn{You manage PiLens over the wire, never over its own Wi-Fi.}
  \item The Pi's built-in Wi-Fi (this becomes the test network).
  \item A phone or laptop you own, to be the ``device under test''.
\end{itemize}

\subsection{The steps}
\begin{enumerate}[leftmargin=1.6em, itemsep=2pt]
  \item Copy the release onto the Pi and run \key{sudo ./install.sh}. It explains
    everything before it touches anything, backs up your config, and \textbf{rolls
    back if anything fails}.
  \item It prints three things: a \textbf{web address}, a \textbf{certificate
    fingerprint}, and a one-time \textbf{setup token}. Keep them.
  \item Open the web address from a computer on your wired network. Your browser
    will warn about the certificate --- that is expected (see \S2). Check the
    fingerprint matches, then continue.
  \item Work through the \textbf{setup wizard}: create your login, choose the
    management network, name the Wi-Fi, pick a password, done.
  \item Connect your test device to the new Wi-Fi network.
  \item Open \textbf{Live Capture}, press \textbf{Start capture} on the Wi-Fi
    interface, and watch the packets roll in.
\end{enumerate}

\begin{heads}
The moment you create the access point, the Pi's own Wi-Fi \emph{becomes} that
access point. If you were connected to the Pi over Wi-Fi you will be kicked off.
\textbf{Always manage it over Ethernet.}
\end{heads}

\screen{overview}{The Overview dashboard --- appliance status, forwarding state, clients, and a live glance at everything at once.}

% ======================= 2. WHAT IT IS =======================
\section{What PiLens is (and isn't)}
\begin{tldr}
PiLens is a \textbf{router with X-ray vision}. Devices on its Wi-Fi go through it
to reach the internet, and it records every packet. It shows you \emph{metadata and
packets} --- it does \textbf{not} decrypt HTTPS.
\end{tldr}

PiLens sits between a test device and the internet. The device joins its Wi-Fi,
gets an address, and routes out through the Pi's Ethernet. Because \emph{every}
packet crosses the Pi, PiLens can capture and analyse all of it --- addresses,
ports, DNS, TLS handshakes, timing, the lot.

\begin{truth}
\textbf{Encrypted means encrypted.} HTTPS, QUIC, VPNs and encrypted DNS stay
scrambled. PiLens shows you \emph{who} talked to \emph{whom}, \emph{when}, over
\emph{what} --- the server name, the TLS version, the certificate when it is sent
in the clear --- but \textbf{not the contents}. That is TLS doing its job, not
PiLens failing. Decryption needs a separate, explicit step on each device (Mode B, \S3.9), and \textbf{a trusted certificate can never be pushed down the
network} --- the device's owner must install it by hand. Anything that claimed
otherwise would be lying to you.
\end{truth}

\subsection{The two padlocks}
\begin{itemize}[leftmargin=1.4em, itemsep=1pt]
  \item \textbf{The management padlock} is the HTTPS on the web interface. It is
    self-signed, so your browser warns once. Check the fingerprint the installer
    printed, trust it, move on.
  \item \textbf{The interception padlock} (decrypting a device's HTTPS) is a
    \emph{separate}, opt-in, per-device thing: a toggle on the Certificates page.
    See \S3.9.
\end{itemize}

% ======================= 3. THE SCREENS =======================
\section{A tour of the screens}
\begin{tldr}
Left-hand menu, top to bottom: \textbf{Overview} (glance), \textbf{Live Capture}
(the main event), \textbf{Conversations / DNS / TLS / Statistics / Findings}
(ways to slice what you caught), then \textbf{Library, Devices, Network, Health,
Settings} (housekeeping).
\end{tldr}

\subsection{Live Capture --- the main event}
\begin{tldr}
Three stacked panes, exactly like Wireshark: the \textbf{packet list} up top, the
\textbf{protocol tree} in the middle, the \textbf{raw bytes} at the bottom. Click a
packet; the other two follow.
\end{tldr}

\screen{analyser}{Live Capture. Packet list, protocol tree and a hex/ASCII inspector. The status bar along the bottom keeps you honest.}

\textbf{The packet list} streams live and is colour-coded by protocol. It stays
smooth even on a busy network because only the visible rows are drawn.

\textbf{The protocol tree} shows the selected packet decoded field by field, by
TShark (the Wireshark engine). Click a field and the matching bytes light up
below --- but \emph{only} when PiLens can prove those bytes really are that field.
It will never make up a highlight.

\textbf{The byte pane} is a proper hex inspector: offsets, hex, ASCII, search,
copy-as-hex or text.

\begin{tip}
The bottom status bar separates things people usually confuse: \textbf{Captured}
(packets stored), \textbf{Ingest lag} (stored but not yet summarised),
\textbf{UI updates dropped} (your browser fell behind --- harmless), and
\textbf{Capture drops} (the appliance genuinely missed packets --- not harmless).
A dropped \emph{screen update} is not a dropped \emph{packet}.
\end{tip}

\begin{heads}
\textbf{Pause view} and \textbf{Pause capture} are different buttons. ``Pause view''
just freezes the scrolling so you can read; capture keeps running. ``Pause
capture'' actually stops storing packets. ``Stop'' ends the session for good.
\end{heads}

\subsection{Filters --- find the needle}
\begin{tldr}
Type a \textbf{Wireshark display filter} in the bar (e.g. \key{dns},
\key{tcp.port == 443}, \key{ip.addr == 192.168.50.20}) and press Enter. It is
evaluated by TShark, so real Wireshark syntax works.
\end{tldr}

The box turns green when the filter is valid, red when it isn't, with the reason.
Press \key{/} anywhere to jump to it. Right-click a packet or a field to build a
filter without typing.

\subsection{Conversations \& Follow TCP Stream}
\begin{tldr}
\textbf{Conversations} groups traffic by endpoint pair with bytes each way. Hit
\textbf{Follow} on a TCP row to reassemble the whole back-and-forth.
\end{tldr}

\screen{conversations}{Conversations --- who spoke to whom, how much each way, sortable, with one-click Follow and filtering.}

\begin{truth}
``Follow'' shows the \emph{transport stream} --- the bytes as they crossed the wire.
For an HTTPS connection that is \textbf{ciphertext}, and PiLens labels it as such.
It is honest about the difference between reassembled packets and decrypted
content.
\end{truth}

\subsection{DNS Explorer}
\begin{tldr}
Every cleartext name lookup: who asked, for what, what came back, and whether it
failed. Group by name, or filter to errors only.
\end{tldr}

\screen{dns}{DNS Explorer, grouped by name --- queries, answers, record types and result codes.}

\begin{heads}
If a device uses \textbf{encrypted DNS} (DoH/DoT), its lookups are inside TLS and
\emph{won't} appear here --- they look like HTTPS to a resolver. PiLens never
pretends encrypted DNS is visible.
\end{heads}

\subsection{TLS Explorer}
\begin{tldr}
Every TLS/QUIC handshake: server name (SNI), negotiated version, cipher, ALPN,
and any certificate sent in the clear. Metadata only --- marked ``encrypted
pass-through'' on every row.
\end{tldr}

\screen{tls}{TLS Explorer. Real server names and versions, honestly labelled. The certificate tab is usually empty on modern servers --- TLS 1.3 encrypts the certificate.}

\begin{tip}
The version column reports the \emph{true negotiated} version, not the
compatibility value on the wire. QUIC shows as ``TLS 1.3 (QUIC)''. If you ever see
everything claiming ``TLS 1.2'', that's a bug in a lesser tool --- here it's read
from the handshake properly.
\end{tip}

\subsection{Statistics}
\begin{tldr}
The shape of the capture: traffic over time, protocol mix, top talkers, packet
sizes. Toggle bytes vs packets. Every number is a real count, defined in the
panel on the right.
\end{tldr}

\screen{statistics}{Statistics --- a traffic histogram, protocol hierarchy, top talkers and size distribution, with a definitions panel so nothing is a mystery.}

\subsection{Findings}
\begin{tldr}
Observable hiccups --- retransmissions, resets, SERVFAILs, NXDOMAINs, dropped
packets --- each with evidence you can click straight to.
\end{tldr}

\screen{findings}{Findings. Each card is a fact with example packets, a reproducing filter, and a note on benign causes.}

\begin{truth}
Findings is a \textbf{magnifying glass, not a smoke alarm}. A finding means a
pattern was \emph{observed}, not that something is wrong or malicious. Each card
spells out the innocent explanations. PiLens will never tell you ``you've been
hacked'' from a retransmission.
\end{truth}

\subsection{HTTPS inspection (Mode B) --- the opt-in one}
\begin{tldr}
A toggle on \textbf{Certificates}. Download the appliance's CA, install and trust
it on a device you own, flip the switch, and that device's HTTPS shows up decrypted
in the \textbf{HTTP Explorer}. Off by default.
\end{tldr}

\screen{certificates}{Certificates \& HTTPS Inspection --- the toggle, the CA to download, per-device install steps, and the privacy options.}

\begin{truth}
This is the only part that can see inside HTTPS, and it only works \emph{with
consent}: you must install the CA on each device by hand and turn trust on. It
\textbf{cannot} be pushed over the network or via DHCP. Pinned apps, QUIC and
private trust stores still can't be read --- they show up as honest errors, never
silently bypassed. The CA's private key never leaves the Pi and is not
downloadable.
\end{truth}

\textbf{To use it:} Certificates \textgreater\ \emph{Download CA certificate}
\textgreater\ install it on your test device (the page has per-OS steps; iOS needs
the extra ``Certificate Trust Settings'' toggle) \textgreater\ verify the
fingerprint matches \textgreater\ flip \emph{Interception} on.

\begin{heads}
Turning interception on while a device has \emph{not} installed the CA will break
that device's HTTPS (it correctly refuses the unknown certificate). Install first,
toggle second --- and toggle it back off when you're done.
\end{heads}

\begin{tip}
\textbf{Bodies are off by default.} The Explorer shows headers, status and timing
for everything; request/response bodies are only stored if you opt in on the
Certificates page, for a named investigation. Clear them when done.
\end{tip}

\screen{http}{HTTP Explorer --- decrypted transactions from trusting devices, with headers and (opt-in) body viewers.}

\subsection{The housekeeping screens}
\begin{description}[leftmargin=0pt, style=nextline, itemsep=3pt]
  \item[\textcolor{amber}{Capture Library}] Every saved session. Rename, tag,
    add notes, mark \emph{preserved} (never auto-deleted), export, or check
    integrity (SHA-256 per file + a manifest).
  \item[\textcolor{amber}{Devices}] Who has been on the Wi-Fi. Label them, add
    notes, jump to their packets. \mn{Randomised MACs show a badge; PiLens never pretends they identify hardware.}
  \item[\textcolor{amber}{Certificates}] The HTTPS-inspection toggle, CA
    download and device install guide (\S3.9).
  \item[\textcolor{amber}{Network Settings}] Change the access point, DHCP, DNS,
    forwarding and management access. Applied as one transaction with automatic
    rollback.
  \item[\textcolor{amber}{System Health}] CPU, memory, temperature, services,
    storage, interfaces --- each with what was actually measured.
  \item[\textcolor{amber}{Settings}] Theme, density, password, two-factor,
    storage limits, colour rules, and the audit log.
\end{description}

\screen{health}{System Health --- every service reports what it actually did, not just whether a process exists.}

% ======================= 4. RECIPES =======================
\section{Common tasks, as recipes}
\begin{tldr}
Short how-tos for the things you'll actually do. Find the one you need, ignore
the rest.
\end{tldr}

\subsection{``See everything one device does''}
\begin{enumerate}[leftmargin=1.6em, itemsep=1pt]
  \item \textbf{Devices} \textgreater\ find it \textgreater\ \textbf{Packets}. (Or
    label it first so you recognise it.)
  \item You land in Live Capture filtered to that device. Add \key{\&\& dns} or
    \key{\&\& tls} to narrow further.
\end{enumerate}

\subsection{``Save this capture as evidence''}
\begin{enumerate}[leftmargin=1.6em, itemsep=1pt]
  \item Stop the capture. \textbf{Library} \textgreater\ the session \textgreater\
    \textbf{Edit} \textgreater\ tick \textbf{Preserved}, add notes.
  \item \textbf{Integrity} \textgreater\ \textbf{Verify hashes}, then \textbf{Download
    manifest}. The original \texttt{.pcapng} opens in desktop Wireshark too.
\end{enumerate}

\subsection{``Analyse a capture from somewhere else''}
\begin{enumerate}[leftmargin=1.6em, itemsep=1pt]
  \item \textbf{Library} \textgreater\ \textbf{Import pcap/pcapng}. It indexes in
    the background; the file is kept exactly as you uploaded it.
\end{enumerate}

\subsection{``Stop a device reaching the internet but still watch it''}
\begin{enumerate}[leftmargin=1.6em, itemsep=1pt]
  \item \textbf{Network Settings} \textgreater\ \textbf{Forwarding} \textgreater\
    \textbf{No forwarding} \textgreater\ apply. Clients still get an address and
    DNS, so you can see what they \emph{try} to reach.
\end{enumerate}

% ======================= 5. TROUBLE =======================
\section{When something's off}
\begin{tldr}
Run \key{sudo /opt/pilens/diagnose.sh} first (it redacts secrets). Then the table.
\end{tldr}

\begin{tabular}{@{}p{5.2cm}p{10.5cm}@{}}
\textbf{Symptom} & \textbf{Try this}\\[-0.5ex]
\multicolumn{2}{@{}l}{\textcolor{line}{\rule{\linewidth}{0.4pt}}}\\
Browser cert warning & Expected. Match the fingerprint, trust it.\\
Can't log in & \key{sudo -u pilens /opt/pilens/bin/pilens reset-password -user admin}\\
Locked out by the subnet & On the console: \key{sudo nft delete table inet pilens}, fix it in the UI, re-apply.\\
No internet for clients & Check the Overview \textbf{Forwarding} card and your Ethernet uplink.\\
AP won't start & Try WPA2 / 2.4\,GHz / a legal channel. Failed applies roll back by themselves.\\
Capture shows drops & Use a capture filter; record to an SSD, not the SD card.\\
Everything's wonky & \key{sudo /opt/pilens/repair.sh} fixes permissions and restarts services.\\
\end{tabular}

\begin{heads}
If a network change ever strands you, every apply keeps a backup and rolls back
on failure, and the last-good firewall is restored at boot before the Wi-Fi comes
up. You are rarely more than a reboot from a working appliance.
\end{heads}

% ======================= 6. CHEAT SHEET =======================
\section{Keyboard cheat sheet}
\begin{center}
\begin{tabular}{@{}ll@{\hspace{10mm}}ll@{}}
\key{/} & focus the filter & \key{Space} & pause / resume the view\\
\key{$\uparrow$}\,\key{$\downarrow$} & move through packets & \key{f} & follow newest packet\\
\key{PgUp}\,\key{PgDn} & jump a page & \key{Shift+F} & focus mode (hide chrome)\\
\key{$\leftarrow$}\,\key{$\rightarrow$} & collapse / expand tree & \key{Alt+1..7} & jump between pages\\
\key{Home}\,\key{End} & first / last packet & right-click & build a filter\\
\end{tabular}
\end{center}

\vspace{2ex}
\begin{truth}
\textbf{The one rule that matters:} only point PiLens at devices and networks you
own or are explicitly allowed to inspect. It has no hidden, covert or
impersonation features, on purpose. Captured traffic can hold passwords and
private messages --- treat your captures like the sensitive data they are, and
delete what you don't need. Access and exports are logged.
\end{truth}

\vfill
\begin{center}\small\color{muted}
PiLens --- built to be fast, honest, and genuinely useful.\\
Full docs live in the repository's \texttt{docs/} folder.
\end{center}

\end{document}
